Los sistemas robóticos y de VR/AR modernos operan sobre un espectro amplio de plataformas, desde nano drones y robots móviles con recursos estrictamente acotados hasta headsets de alta gama. En todos los casos, la eficiencia del sistema de percepción es una propiedad fundamental: en las plataformas más restringidas es una condición de viabilidad, y en las más capaces libera recursos críticos para otras tareas como el renderizado o la planificación de movimiento. En este contexto, la localización precisa y robusta es un requisito ineludible, y el VI-SLAM geométrico, que fusiona cámaras e IMU sin depender de hardware especializado, representa el enfoque más prometedor para satisfacerlo de manera eficiente. Sin embargo, los sistemas del estado del arte aún presentan dificultades para garantizar simultáneamente precisión, robustez y operación en tiempo real en condiciones realistas y trayectorias de larga duración. Este trabajo surge de la necesidad de desarrollar un sistema VI-SLAM que sea a la vez preciso, eficiente y robusto, y que pueda desplegarse en la práctica sobre el amplio espectro de plataformas reales.
